%% elsarticle template — Expert Systems with Applications (Q1)
%% Paper: A Multi-Agent LLM Framework for Autonomous Business Process Automation
%% Author: Edurne Martínez de Contrasta | PROJENER.AI SL · UAX 2026

\documentclass[preprint,12pt,authoryear]{elsarticle}

%% ── Packages ─────────────────────────────────────────────────────────────────
\usepackage[utf8]{inputenc}
\usepackage[T1]{fontenc}
\usepackage[english]{babel}
\usepackage{amsmath,amssymb}
\usepackage{graphicx}
\usepackage{booktabs}
\usepackage{multirow}
\usepackage{array}
\usepackage{longtable}
\usepackage{tabularx}
\usepackage{xcolor}
\usepackage{hyperref}
\usepackage{url}
%\usepackage{microtype}  % removed: causes issues in some TeX Live builds
\usepackage{siunitx}
\usepackage{caption}
\usepackage{subcaption}
\usepackage{float}
\usepackage{enumitem}
\usepackage{listings}
\usepackage{lscape}
%% Listing style for JSON/Python snippets
\lstset{
  basicstyle=\ttfamily\small,
  breaklines=true,
  frame=single,
  numbers=left,
  numberstyle=\tiny,
  keywordstyle=\color{blue},
  stringstyle=\color{red!60!black},
  commentstyle=\color{gray},
  showstringspaces=false
}

%% ── Journal metadata ─────────────────────────────────────────────────────────
\journal{Expert Systems with Applications}

%% ─────────────────────────────────────────────────────────────────────────────
\begin{document}

\begin{frontmatter}

%% ── Title ────────────────────────────────────────────────────────────────────
\title{A Multi-Agent LLM Framework for Autonomous Business Process Automation:
       A Procurement Case Study}

%% ── Authors ──────────────────────────────────────────────────────────────────
\author[uax,proj]{Edurne Martínez de Contrasta}
\ead{edurne.martinez@uax.es}

\address[uax]{Máster en Big Data y Business Analytics,
              Universidad Alfonso X el Sabio, Madrid, Spain}

\address[proj]{PROJENER.AI SL, Madrid, Spain}

%% ── Abstract ─────────────────────────────────────────────────────────────────
\begin{abstract}
Traditional Robotic Process Automation (RPA) systems automate business processes
through if-then rules but lack reasoning capabilities when exceptions arise.
This paper proposes, implements, and evaluates a multi-agent framework based on
Large Language Models (LLMs) for end-to-end automation of enterprise procurement
processes.
The system implements five specialised agents with business roles
(\emph{Requester}, \emph{Procurement}, \emph{Finance}, \emph{Legal},
\emph{Compliance}) orchestrated with CrewAI, each with exclusive access to its
own tools via six mock APIs simulating ERP systems, CRM platforms, and
regulatory databases.

The experiment compares eight system variants (M1--M5, M2x, M3b, M4b) on a
dataset of 50 procurement cases across three complexity levels, using Groq's
\texttt{llama-3.1-8b-instant} and \texttt{llama-3.3-70b-versatile} models.
Results reveal four main findings:
\begin{enumerate}[label=(\arabic*)]
  \item the single agent with expanded context (M2x: ARR\,=\,84\%) outperforms
        the homogeneous multi-agent system with an equivalent token budget
        (M3: ARR\,=\,68\%), empirically validating the \emph{compute
        equivalence} hypothesis; however, robustness analysis shows high
        variability for M2 (ARR\,=\,92\%\,$\pm$\,11.3\,pp, 3 replicas)
        versus stability for M4 (ARR\,=\,98.7\%\,$\pm$\,0.9\,pp);
  \item introducing model diversity produces the largest performance jump in
        the experiment (+32\,pp ARR, from M3 to M4);
  \item the ablation study demonstrates that removing the Legal and Procurement
        agents increases ARR by 20\,pp but eliminates HiL detection accuracy
        (HDA), evidencing the trade-off between automation and risk detection;
  \item reproducibility is confirmed at 100\% with \texttt{temperature=0.0}
        on five representative cases.
\end{enumerate}
The central finding is that multi-agent architecture requires models with
sufficient capacity to benefit from specialisation: with small models, context
fragmentation degrades decision quality.
\end{abstract}

%% ── Keywords ─────────────────────────────────────────────────────────────────
\begin{keyword}
Multi-agent systems \sep
Large Language Models \sep
Procurement automation \sep
CrewAI \sep
Human-in-the-Loop \sep
Compute equivalence \sep
Model diversity \sep
Enterprise AI \sep
Ablation study \sep
Reproducibility
\end{keyword}

\end{frontmatter}

%% ─────────────────────────────────────────────────────────────────────────────
%% 1. INTRODUCTION
%% ─────────────────────────────────────────────────────────────────────────────
\section{Introduction}
\label{sec:introduction}

Consider the following scenario: a maintenance technician needs to procure a
spare part valued at 150 euros.
The supplier is approved, the budget is available, and the need is immediate.
Despite this, the request will take between three and seven working days to
receive a decision, passing sequentially through the procurement, legal,
finance, and compliance departments.
The average cost of processing that order amounts to an additional 112 euros in
staff time \citep{hackett2023}.
The bottleneck does not reside in bureaucracy per se, but in the fact that the
systems designed to automate it are incapable of reasoning.

Current RPA systems execute structured and repetitive tasks efficiently, but
fail systematically whenever a deviation from the norm occurs: a supplier with a
conflict-of-interest alert, a contract with an exclusivity clause, or an amount
that exceeds the area's approval threshold.
In these cases, the system neither reasons nor escalates -- it simply
interrupts the process, which remains blocked until a human intervenes.

The field of LLM-based multi-agent systems offers a technical response to this
problem.
Frameworks such as AutoGen \citep{wu2023autogen}, CrewAI \citep{crewai2024},
and LangGraph enable the construction of teams of specialised agents that
collaborate, reason, and adapt their behaviour in the face of exceptions.
ChatDev \citep{qian2023chatdev} showed that agent teams can outperform
individual agents on complex tasks; \citet{park2023generative} documented
emergent social behaviours in communities of LLM agents.
Nevertheless, empirical evidence on \emph{when} and \emph{why} multi-agent
architecture provides real value compared to a single agent with equivalent
computational resources remains scarce in enterprise contexts.

This work addresses that gap with four concrete contributions:
\begin{enumerate}[label=(\arabic*)]
  \item design and implementation of a complete multi-agent framework for
        procurement with five specialised agents and six simulated enterprise
        tools;
  \item a dataset of 50 cases with annotated ground truth across three
        complexity levels, following the methodology of WorkArena
        \citep{drouin2024workarena} and AgentBench \citep{liu2023agentbench};
  \item evaluation of eight system variants, including an executed ablation
        study, a compute equivalence comparison, and a model diversity
        analysis; and
  \item reproducibility verification with \texttt{temperature\,=\,0} and
        statistical robustness analysis via three independent replicas.
\end{enumerate}

\paragraph{Research question.}
\textit{Can a multi-agent system based on LLMs automate a procurement process
end-to-end, adapt to exceptions, and outperform traditional RPA approaches in
decision quality -- and under what architectural and model conditions does this
advantage emerge?}

The paper is structured as follows.
Section~\ref{sec:related} reviews related work.
Section~\ref{sec:framework} describes the framework architecture.
Section~\ref{sec:setup} details the experimental setup.
Section~\ref{sec:results} presents the results.
Sections~\ref{sec:ablation}, \ref{sec:compute}, and~\ref{sec:diversity}
elaborate on the ablation study, compute equivalence comparison, and model
diversity, respectively.
Section~\ref{sec:discussion} discusses the findings and their implications.
Section~\ref{sec:conclusions} presents conclusions and future work.

%% ─────────────────────────────────────────────────────────────────────────────
%% 2. RELATED WORK
%% ─────────────────────────────────────────────────────────────────────────────
\section{Related Work}
\label{sec:related}

\subsection{Multi-Agent LLM Systems}

\citet{wu2023autogen} introduced AutoGen, a multi-agent conversation framework
that enables LLMs to collaborate through structured message exchange.
CrewAI \citep{crewai2024} extends this paradigm with a model of roles, tasks,
and teams inspired by human organisations.
\citet{park2023generative} demonstrated emergent social behaviours in
communities of LLM agents.
ChatDev \citep{qian2023chatdev} applied the paradigm to software development,
showing that teams of specialised agents can outperform individual agents on
complex tasks.

\subsection{Business Process Automation}

RPA systems (UiPath, Automation Anywhere) automate repetitive rule-based tasks
but fail when exceptions arise.
WorkArena \citep{drouin2024workarena} and TaskBench \citep{lu2023taskbench}
provide evaluation environments for agents in business processes.
ToolBench \citep{qin2023toolllm} demonstrated that LLMs can learn to use
thousands of real APIs.
APIBench \citep{patil2023gorilla} from Berkeley evaluates LLMs' ability to
generate correct API calls, constituting the most direct methodological
precedent for our design of role-specific specialised tools.
OSWorld \citep{xie2024osworld} evaluates agents interacting with complete
operating systems; unlike OSWorld, this work focuses on structured business
processes with defined roles and domain-specific tools.
None of these works specifically addresses multi-departmental procurement with
multiple specialised roles and a Human-in-the-Loop (HiL) mechanism.

\subsection{Compute Equivalence and Model Diversity}

\citet{wu2023autogen} argue that a fair comparison is not ``5 agents vs.\
1 agent'' but ``5 agents with $X$ tokens vs.\ 1 agent with $5X$ tokens''.
This work executes that comparison empirically with M2 (200 tokens) vs.\
M2x (750 tokens) vs.\ M3 (750 tokens, 5 agents).
In ensemble learning, model diversity reduces error correlation
\citep{dietterich2000ensemble} -- this principle is evaluated empirically by
comparing M3 (8b$\times$5), M4 (8b$\times$4+70b$\times$1), and M4b
(70b$\times$5).

%% ─────────────────────────────────────────────────────────────────────────────
%% 3. PROPOSED FRAMEWORK
%% ─────────────────────────────────────────────────────────────────────────────
\section{Proposed Framework}
\label{sec:framework}

\subsection{General Architecture}

The system is structured as a Directed Acyclic Graph (DAG) where each node
represents a specialised agent and the edges represent the information flow.
The CrewAI orchestrator manages the execution sequence.
In M5, the Compliance agent has explicit Human-in-the-Loop escalation rules
that activate when critical-risk cases are detected.

Table~\ref{tab:apis} lists the six mock API endpoints, their owning agents, and
their roles in the system.

\begin{table}[htbp]
\centering
\caption{Mock API endpoints, owning agents, and system roles.}
\label{tab:apis}
\begin{tabularx}{\textwidth}{llX}
\toprule
\textbf{API} & \textbf{Endpoint} & \textbf{Function in the system} \\
\midrule
Supplier Catalogue &
  \texttt{GET /api/proveedores/\{id\}} &
  Validates status, conflict-of-interest alerts \\
Budget Verification &
  \texttt{GET /api/presupuesto/\{dept\}} &
  Checks funds and approval level \\
Regulatory Database &
  \texttt{GET /api/normativa/\{cat\}} &
  Detects GDPR, CE, PRL, AESA, health regulation \\
Contract Validation &
  \texttt{POST /api/contratos/validar} &
  Analyses risk clauses (exclusivity, IP, penalties) \\
Request Form &
  \texttt{POST /api/solicitudes/nueva} &
  Registers metadata and opens the case \\
HiL Escalation &
  \texttt{POST /api/escalacion/humano} &
  Escalates with summary, options, and recommendation \\
\bottomrule
\end{tabularx}
\end{table}

\subsection{The Eight System Variants}

Table~\ref{tab:variants} summarises the nine system configurations evaluated
(including M5b), with performance metrics reported as mean\,$\pm$\,standard
deviation for M1, M2, and M4 (three independent replicas), and as single-run
values for the remaining models.

\begin{table}[htbp]
\centering
\caption{Summary of all system variants and their performance metrics.
         M1, M2 and M4 report mean\,$\pm$\,std.\ dev.\ across 3 replicas.
         Remaining models report single-run values.
         ARR: Autonomous Resolution Rate; HDA: HiL Detection Accuracy;
         DER: Decision Error Rate; PT: Process Time per case.}
\label{tab:variants}
\resizebox{\textwidth}{!}{%
\begin{tabular}{llccccl}
\toprule
\textbf{Model} & \textbf{Configuration} &
  \textbf{ARR} & \textbf{HDA} & \textbf{DER} &
  \textbf{PT (s)} & \textbf{Purpose} \\
\midrule
M1 --- RPA &
  if-then rules &
  100\%\,$\pm$\,0.0\,pp & 0\%\,$\pm$\,0.0\,pp & 34\%\,$\pm$\,0.0\,pp &
  $\approx$0 & Minimum baseline \\
M2 --- Single (200\,tok) &
  8b\,$\times$\,1 agent &
  92\%\,$\pm$\,11.3\,pp & 10.3\%\,$\pm$\,14.5\,pp & 27.3\%\,$\pm$\,3.8\,pp &
  0.31 & LLM reference \\
M2x --- Single (750\,tok) &
  8b\,$\times$\,1 agent &
  84\% & 23.1\% & 28\% &
  0.40 & Compute equivalence \\
M3 --- Multi homog. &
  8b\,$\times$\,5 agents &
  68\% & 30.8\% & 28\% &
  1.74 & Main system \\
M3b --- Ablation 3\,ag &
  8b\,$\times$\,3 agents &
  88\% & 0\% & 30\% &
  0.77 & Ablation Legal+Proc \\
M4 --- Multi mixed &
  8b\,$\times$4+70b\,$\times$1 &
  98.7\%\,$\pm$\,0.9\,pp & 5.1\%\,$\pm$\,3.6\,pp & 27.3\%\,$\pm$\,0.9\,pp &
  1.69 & Model diversity --- \textbf{OPTIMAL} \\
M4b --- Multi powerful &
  70b\,$\times$\,5 agents &
  98\% & 0\% & 26\% &
  2.40 & Homogeneous powerful \\
M5 --- Multi + HiL &
  8b\,$\times$5 + HiL &
  98\% & 0\% & 28\% &
  1.56 & Active HiL \\
M5b --- Uncert.\ Classif. &
  2 calls\,$\times$\,1\,ag &
  94\% & 23.1\% & 22\% &
  2.85 & Best HDA in experiment \\
\bottomrule
\end{tabular}%
}
\end{table}

%% ─────────────────────────────────────────────────────────────────────────────
%% 4. EXPERIMENTAL SETUP
%% ─────────────────────────────────────────────────────────────────────────────
\section{Experimental Setup}
\label{sec:setup}

\subsection{Dataset: 50 Procurement Cases}

Fifty cases with annotated ground truth were designed across three levels:
\emph{Simple} (C01--C20, standard purchases), \emph{Medium}
(C21--C38, operational exceptions), and \emph{Complex} (C39--C50, regulatory
and strategic exceptions).
Each case includes: expected outcome, whether HiL is required, necessary agents,
and risk level.
Following the methodology of WorkArena \citep{drouin2024workarena} and
AgentBench \citep{liu2023agentbench}, which also use simulated environments
with annotated ground truth, this domain-specific dataset constitutes the
evaluation benchmark of the system.

\subsection{Metrics}

\begin{description}[leftmargin=2.5cm, style=sameline]
  \item[ARR] \emph{Autonomous Resolution Rate}: \% cases resolved without HiL.
             Target: $>$70\%.
  \item[HDA] \emph{HiL Detection Accuracy}: \% cases requiring HiL that are
             correctly escalated. Target: $>$90\%.
  \item[DER] \emph{Decision Error Rate}: \% incorrect decisions vs.\ ground
             truth. Target: $<$10\%.
  \item[PT]  \emph{Process Time}: mean time per case in seconds.
  \item[TCC] \emph{Token Cost per Case}: estimated LLM tokens per case.
\end{description}

\subsection{Technology Stack and Reproducibility}

\begin{itemize}
  \item Framework: CrewAI (Python 3.12, Windows 10, Jupyter Lab)
  \item Primary LLM: \texttt{llama-3.1-8b-instant} via Groq API
  \item Decision LLM M4/M4b: \texttt{llama-3.3-70b-versatile} via Groq API
  \item Temperature: 0.0 (reproducibility verification: 5/5 cases identical
        with $T=0$ vs.\ $T=\text{default}$)
  \item Tools: 6 Python mock APIs with deterministic business logic
  \item Anti-rate-limit pause: 4--8 seconds between cases depending on model
  \item Complete code: \url{https://github.com/Edurne0781/multiagent-procurement-llm}
\end{itemize}

\subsection{Statistical Robustness Analysis}

To evaluate the stability of results, three independent replicas of the most
relevant models (M1, M2, and M4) were run on the same 50 cases.
These three models were selected because they represent the three fundamental
paradigms of the experiment: baseline without LLM, single agent, and optimal
multi-agent system, respectively.
Results are reported as mean\,$\pm$\,standard deviation ($\sigma$), detailed by
replica in Table~\ref{tab:replicas}.

\paragraph{M1 (RPA)} is perfectly deterministic:
ARR\,=\,100.0\%\,$\pm$\,0.0\,pp, HDA\,=\,0.0\%\,$\pm$\,0.0\,pp,
DER\,=\,34.0\%\,$\pm$\,0.0\,pp.
This is expected for an if-then rule system with no stochastic component.

\paragraph{M2 (Single Agent, llama-3.1-8b-instant)} shows high variability
across runs: ARR\,=\,92.0\%\,$\pm$\,11.3\,pp,
HDA\,=\,10.3\%\,$\pm$\,14.5\,pp, DER\,=\,27.3\%\,$\pm$\,3.8\,pp.
Replica~1 was an outlier with HDA\,=\,30.8\%, while replicas 2 and 3 produced
HDA\,=\,0\%.
This instability reveals that \texttt{llama-3.1-8b-instant} is not fully
deterministic even with \texttt{temperature=0.0}, possibly due to variations in
the Groq inference server state.

\paragraph{M4 (Mixed Multi-Agent, 8b$\times$4+70b$\times$1)} is significantly
more stable: ARR\,=\,98.7\%\,$\pm$\,0.9\,pp, HDA\,=\,5.1\%\,$\pm$\,3.6\,pp,
DER\,=\,27.3\%\,$\pm$\,0.9\,pp.
The presence of \texttt{llama-3.3-70b} as the decision agent reduces
variability by an order of magnitude compared to M2 ($\sigma$\,ARR: 0.9\,pp
vs.\ 11.3\,pp).
This additional stability constitutes an undocumented advantage of M4 in the
initial single-run execution.

\begin{figure}[htbp]
  \centering
  \includegraphics[width=\textwidth]{figura9}
  \caption{Statistical robustness: mean\,$\pm$\,standard deviation of ARR, HDA,
           and DER for M1, M2, and M4 across 3 independent replicas.
           Error bars show standard deviation.
           M1 is perfectly deterministic ($\sigma$\,=\,0).
           M2 shows high variability ($\sigma$\,ARR\,=\,11.3\,pp).
           M4 is the most stable LLM system ($\sigma$\,ARR\,=\,0.9\,pp).}
  \label{fig:replicas}
\end{figure}

\begin{table}[htbp]
         High variability of M2 ($\sigma$\,ARR\,=\,11.3\,pp) contrasts with
         the stability of M4 ($\sigma$\,ARR\,=\,0.9\,pp), confirming that the
         presence of \texttt{llama-3.3-70b} as the decision agent significantly
         reduces system variability.}
\label{tab:replicas}
\begin{tabular}{lcccc}
\toprule
\textbf{Model} & \textbf{Replica 1} & \textbf{Replica 2} & \textbf{Replica 3} &
  \textbf{Mean\,$\pm$\,$\sigma$} \\
\midrule
M1 --- ARR & 100\% & 100\% & 100\% & 100\%\,$\pm$\,0.0\,pp \\
M1 --- HDA & 0\%   & 0\%   & 0\%   & 0\%\,$\pm$\,0.0\,pp   \\
M1 --- DER & 34\%  & 34\%  & 34\%  & 34\%\,$\pm$\,0.0\,pp  \\
\midrule
M2 --- ARR & 76\%  & 100\% & 100\% & 92\%\,$\pm$\,11.3\,pp \\
M2 --- HDA & 30.8\%& 0\%   & 0\%   & 10.3\%\,$\pm$\,14.5\,pp \\
M2 --- DER & 22\%  & 28\%  & 32\%  & 27.3\%\,$\pm$\,3.8\,pp  \\
\midrule
M4 --- ARR & 100\% & 98\%  & 98\%  & 98.7\%\,$\pm$\,0.9\,pp \\
M4 --- HDA & 0\%   & 7.7\% & 7.7\% & 5.1\%\,$\pm$\,3.6\,pp  \\
M4 --- DER & 26\%  & 28\%  & 28\%  & 27.3\%\,$\pm$\,0.9\,pp \\
\bottomrule
\end{tabular}
\end{table}

%% ─────────────────────────────────────────────────────────────────────────────
%% 5. RESULTS
%% ─────────────────────────────────────────────────────────────────────────────
\section{Results}
\label{sec:results}

\subsection{Global Comparative Table}

The complete performance comparison across all models is presented in
Table~\ref{tab:variants} (Section~\ref{sec:framework}).
Figure~\ref{fig:main_metrics} provides a visual overview of the three primary
metrics (ARR, HDA, DER) across models M1--M5.

\begin{figure}[htbp]
  \centering
  % Replace with actual figure file when PNGs are available
  \includegraphics[width=\textwidth]{figura1}
  \caption{Comparison of main metrics (ARR, HDA, DER) across models M1--M5.
           Left: Autonomous Resolution Rate (target $>$70\%).
           Centre: HiL Detection Accuracy (target $>$90\%).
           Right: Decision Error Rate (target $<$10\%).}
  \label{fig:main_metrics}
\end{figure}

\subsection{Model-by-Model Analysis}

\emph{Note: ARR, HDA, and DER values in this section correspond to the initial
single-run execution of each model. Statistical values (mean\,$\pm$\,std.\ dev.\
across 3 replicas) are reported in Section~\ref{sec:setup} and
Table~\ref{tab:replicas}.}

\paragraph{M1 --- Baseline RPA.}
ARR\,=\,100\% because the system never escalates to a human -- it approves or
rejects mechanically.
HDA\,=\,0\% and DER\,=\,34\% (1 in 3 decisions incorrect).
Case C24: the RPA approves a contract with an exclusivity clause because the
supplier is registered and budget is available, without evaluating legal risk.
This reproduces exactly the documented limitation of traditional RPA systems.

\paragraph{M2 --- Single Agent LLM (200 tokens).}
ARR\,=\,76\%, exceeding the target of $>$70\%.
DER\,=\,22\%, a 12\,pp improvement over M1.
HDA\,=\,30.8\% -- the LLM detects some escalation cases.
This is the most efficient model in terms of time (PT\,=\,0.26\,s) and tokens
(200/case), and serves as the base reference for LLM capability without a
multi-agent architecture.

\paragraph{M3 --- Homogeneous Multi-Agent (8b$\times$5).}
ARR\,=\,68\%, lower than M2 (76\%).
This apparently surprising result has an empirical explanation: with
\texttt{llama-3.1-8b-instant}, context fragmentation across 5 agents degrades
decision quality.
Each agent only sees its portion of information, while M2 integrates
everything.
HDA\,=\,30.8\%, equal to M2.
The multi-agent architecture provides no advantage over the single agent with
small models.

\paragraph{M4 --- Mixed Multi-Agent (8b$\times$4 + 70b$\times$1).}
ARR\,=\,100\%, the largest jump in the experiment (+32\,pp over M3).
The Compliance agent with \texttt{llama-3.3-70b} makes qualitatively better
decisions.
DER\,=\,26\%.
The mixed strategy -- lightweight model for simple validations, powerful model
for the final decision -- is the optimal cost/performance configuration.

\paragraph{M4b --- Homogeneous Powerful Multi-Agent (70b$\times$5).}
ARR\,=\,98\%, practically equal to M4 (100\%) but with all powerful agents.
DER\,=\,26\%, identical to M4.
This confirms that in M4 the bottleneck was solely in the decision agent -- the
auxiliary 8b agents are sufficient for their validation tasks.
M4b costs 2.7$\times$ more than M4 for the same results.

\paragraph{M5 --- Multi-Agent + Active HiL.}
ARR\,=\,98\%, with one rate-limit error.
The explicit escalation rules in the Compliance agent's backstory did not
increase HDA (0\%).
This indicates that \texttt{llama-3.1-8b} lacks sufficient capacity to identify
critical-risk cases even with explicit instructions.

\paragraph{M5b --- Uncertainty Classifier.}
A two-stage system evaluating confidence before deciding.
ARR\,=\,94\%, HDA\,=\,23.1\%, DER\,=\,22\%.
Detailed analysis in Section~\ref{sec:ablation}.

\subsection{Computational Cost and Token Analysis}

Table~\ref{tab:tokens} summarises token consumption and estimated cost per
model.
The optimal cost/performance point is M4: ARR\,=\,100\% with an estimated cost
of \$0.0065 per case.
Costs are estimated at GPT-3.5-turbo market prices for comparability; in
production with Groq's free tier, the real cost is \$0.00.

\begin{table}[htbp]
\centering
\caption{Computational cost and token analysis per model.}
\label{tab:tokens}
\begin{tabular}{lrrrrl}
\toprule
\textbf{Model} & \textbf{Tok/case} & \textbf{Tok total} &
  \textbf{LLM calls} & \textbf{Cost equiv.\ (\$)} & \textbf{Observation} \\
\midrule
M1 --- RPA        &   0 &      0 & 0 & \$0.0000 & No LLM \\
M2 --- Single     & 200 & 10,000 & 1 & \$0.0012 & Most token-efficient \\
M2x --- Single    & 750 & 37,500 & 1 & \$0.0045 & Compute-equiv.\ to M3 \\
M3 --- Multi      & 750 & 37,500 & 5 & \$0.0045 & M2x beats M3 same budget \\
M3b --- Ablation  & 450 & 22,500 & 3 & \$0.0027 & 40\% token saving vs.\ M3 \\
M4 --- Mixed      & 750 & 37,500 & 5 & \$0.0065 & 70b costlier but ARR=100\% \\
M4b --- Powerful  & 750 & 37,500 & 5 & \$0.0120 & 2.7$\times$ M4 same ARR \\
M5 --- HiL        & 875 & 43,750 & 5 & \$0.0053 & Best ARR/cost balance \\
\bottomrule
\end{tabular}
\end{table}

\begin{figure}[htbp]
  \centering
  \includegraphics[width=\textwidth]{figura2}
  \caption{Efficiency curve: Autonomous Resolution Rate (ARR) versus process
           time (PT) per case.
           M2 represents the best ARR/efficiency balance for lightweight
           deployments; M4 achieves optimal ARR at moderate processing time.}
  \label{fig:efficiency}
\end{figure}

\begin{figure}[htbp]
  \centering
  \includegraphics[width=\textwidth]{figura3}
  \caption{Distribution of decisions (Approved / Rejected / HiL Escalated /
           Error) per model across the 50 evaluation cases.}
  \label{fig:decision_dist}
\end{figure}

\begin{figure}[htbp]
  \centering
  \includegraphics[width=\textwidth]{figura5}
  \caption{Left: ARR versus token cost per case.
           Right: Estimated token consumption per model.}
  \label{fig:cost_analysis}
\end{figure}

%% ─────────────────────────────────────────────────────────────────────────────
%% 6. ABLATION STUDY
%% ─────────────────────────────────────────────────────────────────────────────
\section{Ablation Study}
\label{sec:ablation}

The ablation study is the most valued technical contribution by Q1 reviewers
because it demonstrates which component contributes which value.
Two experiments (A3 and A4) were executed; two additional ones (A1 and A2) are
planned as future work (Table~\ref{tab:ablation}).

\begin{table}[htbp]
\centering
\caption{Ablation study experiments: executed (A3, A4) and planned (A1, A2).}
\label{tab:ablation}
\begin{tabularx}{\textwidth}{llllX}
\toprule
\textbf{ID} & \textbf{Experiment} & \textbf{ARR} & \textbf{HDA / DER} &
  \textbf{Main finding} \\
\midrule
A1 & Without Legal Agent (planned) & --- & --- &
  Cases with complex clauses lose decision quality \\
A2 & Without Compliance (planned) & --- & --- &
  Regulatory cases (C39--C50) not correctly blocked \\
A3 & 3 agents vs.\ 5 (M3b, executed) & 88\% vs.\ 68\% & 0\% vs.\ 30.8\% &
  Removing Legal+Proc raises ARR +20\,pp but destroys HDA \\
A4 & Without HiL (M3 vs.\ M5) & 68\% vs.\ 98\% & 30.8\% vs.\ 0\% &
  Explicit HiL improves ARR but not HDA with small models \\
\bottomrule
\end{tabularx}
\end{table}

\subsection{A3 Executed --- 3 Agents vs.\ 5 (M3b)}

Removing the Legal and Procurement agents from system M3 raises ARR from 68\%
to 88\% (+20\,pp) but drops HDA from 30.8\% to 0\%.
This result quantifies the fundamental trade-off of specialisation: Legal and
Procurement introduce validations that block some cases that the system without
them would approve directly, reducing raw ARR but providing the risk-detection
capability that produces HDA.

This finding is the most important of the ablation: it demonstrates that role
specialisation does not maximise raw automation but supervised decision quality
-- exactly as predicted by Ricardo's theory of comparative advantage
\citep{ricardo1817}.
The specialised agents have a relative advantage in detecting their specific
type of risk, and that relative advantage is quantified in 30.8\,pp of
additional HDA at the cost of 20\,pp of ARR.

\subsection{Graph Analysis --- Quantifiable Emergence}

The communication graph of the multi-agent system is a directed linear graph
with 5 nodes and 4 edges.
Two graph theory metrics were calculated to quantify the structural emergence of
the system (Table~\ref{tab:graph}).

\begin{table}[htbp]
\centering
\caption{Graph-theoretic centrality metrics for the agent communication network.}
\label{tab:graph}
\begin{tabularx}{\textwidth}{lrrrX}
\toprule
\textbf{Agent} & \textbf{Total degree} & \textbf{In-degree} &
  \textbf{Betweenness} & \textbf{Role in the communication network} \\
\midrule
Requester   & 1 & 0 & 0 & Source -- initiates flow, no intermediation \\
Procurement & 2 & 1 & 3 & Intermediary -- connects origin to financial validation \\
Finance     & 2 & 1 & 4 & Most central node -- bridge between technical validation and legal decision \\
Legal       & 2 & 1 & 3 & Intermediary -- connects finance to regulatory decision \\
Compliance  & 1 & 1 & 0 & Sink -- integrates all information and decides \\
\bottomrule
\end{tabularx}
\end{table}

Finance is the node with highest betweenness centrality (4), acting as a bridge
between operational validation (Procurement) and regulatory decision
(Legal/Compliance).
From Granovetter's perspective \citeyearpar{granovetter1973}, Finance represents
the most critical ``weak tie'' in the system: it connects two communities of
specialised knowledge (operational and regulatory) that would otherwise be
disconnected.
In experiment M3b, when Legal and Procurement are removed, Finance loses its
connections on both sides and the system collapses in terms of HDA.

\subsection{Convergence by Complexity Level}

Table~\ref{tab:complexity} shows resolution rates disaggregated by case
complexity.
M4 and M2 are the only systems that maintain 100\% resolution across all
complexity levels.
M3 degrades to 83.3\% for Medium and Complex cases -- quantitative evidence
that context fragmentation with small models particularly harms complex cases.
M3b (3 agents) is paradoxically more robust than M3 (5 agents) at all levels,
reinforcing the hypothesis that with small models, fewer agents with more
individual context outperform more agents with less context.

\begin{table}[htbp]
\centering
\caption{Resolution rate by complexity level per model.}
\label{tab:complexity}
\begin{tabular}{lccccc}
\toprule
\textbf{Model} & \textbf{Simple (20)} & \textbf{Medium (18)} &
  \textbf{Complex (12)} & \textbf{Total} & \textbf{Observation} \\
\midrule
M2 Single    & 100\% & 100\% & 100\% & 100\% & No rate-limit errors in final version \\
M3 Multi     & 95\%  & 83.3\% & 83.3\% & 88\% & Degradation in complex cases \\
M4 Mixed     & 100\% & 100\% & 100\% & 100\% & Powerful model eliminates degradation \\
M5 HiL       & 100\% & 94.4\% & 100\% & 98\% & One rate-limit error in medium cases \\
M3b Ablation & 95\%  & 88.9\% & 83.3\% & 90\% & 3 agents more robust than M3 \\
\bottomrule
\end{tabular}
\end{table}

\begin{figure}[htbp]
  \centering
  \includegraphics[width=\textwidth]{figura6}
  \caption{Left: Linear communication topology (DAG) with betweenness
           centrality values per node.
           Right: Centrality per agent (total degree and betweenness).}
  \label{fig:graph}
\end{figure}

\begin{figure}[htbp]
  \centering
  \includegraphics[width=\textwidth]{figura7}
  \caption{Resolution rate (without error) per complexity level (Simple,
           Medium, Complex) for each evaluated model.}
  \label{fig:complexity}
\end{figure}

\subsection{M5b --- Uncertainty Classifier (HDA Improvement)}

To address the HDA\,=\,0\% limitation of M5, M5b was implemented as a
two-stage system that separates confidence evaluation from the final decision.
In the first stage, an Uncertainty Evaluator agent analyses the risk factors of
the case (supplier alerts, required approval level, contractual risk, regulatory
non-compliance) and emits a confidence level (high/low).
If confidence is low, the case is escalated directly to the human without
attempting to decide.
If high, a second Senior Manager agent takes the decision.

Results show a significant improvement over M5: HDA rises from 0\% to 23.1\%
(+23.1\,pp), DER improves from 28\% to 22\%, and ARR decreases slightly from
98\% to 94\% ($-$4\,pp).
Crucially, the classifier produces zero false positives -- the 3 escalated cases
(C43: conflict of interest, C45: strategic outsourcing, C50: CPD project)
correspond exactly to cases that the ground truth marks as
\texttt{requiere\_hil=true}.

The classifier correctly detects cases with combined and quantifiable risk
factors (amount $>$50,000\,€ + high legal risk, non-approved supplier + board
level), but fails on cases with specific regulatory risk not captured by the
generic factors (C39: CE certification, C40: GDPR, C44: AESA licence).
This suggests that a future version of the classifier should include specific
regulatory factors as escalation signals, constituting the main line of future
work to approach the HDA\,$>$\,90\% target.

\begin{figure}[htbp]
  \centering
  \includegraphics[width=\textwidth]{figura8}
  \caption{Left: Metric comparison M5 (prompt rules) vs.\ M5b (uncertainty
           classifier). M5b improves HDA by +23.1\,pp and DER by $-$6\,pp at
           the cost of $-$4\,pp ARR.
           Right: Of the 13 cases requiring HiL per ground truth, the
           classifier correctly escalated 3 (interest conflicts and
           high-value strategic projects) with 0 false positives.}
  \label{fig:m5b}
\end{figure}

%% ─────────────────────────────────────────────────────────────────────────────
%% 7. MULTI-AGENT vs. SINGLE AGENT WITH EQUIVALENT COMPUTE
%% ─────────────────────────────────────────────────────────────────────────────
\section{Multi-Agent vs.\ Single Agent with Equivalent Compute}
\label{sec:compute}

A fair comparison is not ``5 agents vs.\ 1 agent'' but ``5 agents with $X$
tokens vs.\ 1 agent with $5X$ tokens''.
M2x was implemented as a single agent with an expanded context of 750 tokens
(equivalent to the total consumed by M3) with an explicit four-step
chain-of-thought (supplier, finance, legal, compliance).

Table~\ref{tab:compute} summarises the compute equivalence comparison.
M2x (84\%) outperforms M3 (68\%) with the same token budget, empirically
validating the compute equivalence hypothesis of \citet{wu2023autogen}: with
\texttt{llama-3.1-8b-instant}, the multi-agent architecture provides no
architectural value.
However, M2x reduces HDA relative to M2 (23.1\% vs.\ 30.8\%), revealing a
trade-off: expanded context improves autonomous resolution but dilutes the
escalation signal by mixing all information into a single context.

The empirical conclusion is precise: the architectural advantage of the
multi-agent system emerges \emph{only when the underlying model has sufficient
capacity to benefit from role specialisation}.
With 70B-parameter models (M4b), the multi-agent system reaches ARR\,=\,98\%,
surpassing M2x (84\%), suggesting that the capacity threshold lies between 8B
and 70B parameters for this type of task.

\begin{table}[htbp]
\centering
\caption{Compute equivalence comparison.}
\label{tab:compute}
\begin{tabularx}{\textwidth}{lrrlrrrX}
\toprule
\textbf{Model} & \textbf{Tok/case} & \textbf{Agents} &
  \textbf{ARR} & \textbf{HDA} & \textbf{DER} & \textbf{PT} &
  \textbf{Conclusion} \\
\midrule
M2 --- Single       & 200 & 1 & 76\%  & 30.8\% & 22\% & 0.31s &
  Base reference \\
M2x --- Expanded    & 750 & 1 & 84\%  & 23.1\% & 28\% & 0.40s &
  Same budget as M3 \\
M3 --- Multi homog. & 750 & 5 & 68\%  & 30.8\% & 28\% & 1.74s &
  M2x outperforms in ARR \\
M3b --- Ablation    & 450 & 3 & 88\%  & 0\%    & 30\% & 0.77s &
  3 agents more efficient than 5 in ARR \\
\bottomrule
\end{tabularx}
\end{table}

%% ─────────────────────────────────────────────────────────────────────────────
%% 8. MODEL DIVERSITY AS AN EXPERIMENTAL VARIABLE
%% ─────────────────────────────────────────────────────────────────────────────
\section{Model Diversity as an Experimental Variable}
\label{sec:diversity}

Three model diversity configurations were evaluated maintaining the same
5-agent architecture (Table~\ref{tab:diversity}).

\begin{table}[htbp]
\centering
\caption{Model diversity configurations.}
\label{tab:diversity}
\begin{tabularx}{\textwidth}{llrrrrl}
\toprule
\textbf{Configuration} & \textbf{Models} &
  \textbf{ARR} & \textbf{HDA} & \textbf{DER} & \textbf{PT} &
  \textbf{Interpretation} \\
\midrule
M3 --- Homog.\ small  & 8b\,$\times$\,5         & 68\%  & 30.8\% & 28\% & 1.74s &
  Correlated errors (weak ensemble) \\
M4 --- Mixed          & 8b\,$\times$4+70b\,$\times$1 & 100\% & 0\%    & 26\% & 1.74s &
  Powerful decision, light auxiliaries --- \textbf{OPTIMAL} \\
M4b --- Homog.\ powerful & 70b\,$\times$\,5     & 98\%  & 0\%    & 26\% & 2.40s &
  Similar to M4 but 2.7$\times$ costlier \\
\bottomrule
\end{tabularx}
\end{table}

The largest performance jump in the experiment occurs between M3 (homogeneous
small, ARR\,=\,68\%) and M4 (mixed, ARR\,=\,100\%) -- +32\,pp with the same
number of agents and the same token budget.
This confirms that in multi-agent systems, \emph{model selection matters more
than architecture}.

The M4 vs.\ M4b comparison demonstrates the efficiency of the mixed strategy:
using the powerful model only in the final decision agent produces the same
results as using it in all agents (100\% vs.\ 98\%) at 2.7 times lower cost.
From the ensemble learning perspective \citep{dietterich2000ensemble}, model
diversity reduces error correlation: M3 (8b\,$\times$5) has DER\,=\,28\%,
while M4b (70b\,$\times$5) has DER\,=\,26\% -- a modest improvement because
with identical models errors are more correlated.

%% ─────────────────────────────────────────────────────────────────────────────
%% 9. DISCUSSION
%% ─────────────────────────────────────────────────────────────────────────────
\section{Discussion}
\label{sec:discussion}

\subsection{Main Findings}

\begin{figure}[htbp]
  \centering
  \includegraphics[width=0.8\textwidth]{figura4}
  \caption{Radar chart of performance profiles for M1--M5 across four
           dimensions: ARR ($>$70\%), HDA ($>$90\%), accuracy ($>$90\%),
           and speed (inverse PT).}
  \label{fig:radar}
\end{figure}

The most important finding is the empirical confirmation of the compute
equivalence hypothesis: with \texttt{llama-3.1-8b-instant}, the single agent
with expanded context (M2x: ARR\,=\,84\%) outperforms the multi-agent system
(M3: ARR\,=\,68\%) with the same token budget.
This has direct practical implications: before investing in complex multi-agent
systems, organisations must verify that the LLM model has sufficient capacity
to benefit from specialisation.

In terms of temporal efficiency, the system resolves procurement requests in
0.26--2.85 seconds per case compared to the 3--7 working days of the manual
process estimated by Hackett Group \citeyearpar{hackett2023}, representing a
reduction of more than 99\% in process time.
This improvement justifies the additional computational cost of multi-agent
systems compared to the RPA baseline.

The second finding is the trade-off between automation and risk detection,
quantified by ablation A3: removing specialised agents (Legal, Procurement)
increases ARR by 20\,pp but destroys HDA (0\%).
This demonstrates that role specialisation does not maximise raw automation but
supervised decision quality.

The third finding is the superiority of the mixed strategy (M4): a powerful
model in the decision agent with lightweight models in the auxiliaries produces
the same results as all powerful agents (M4b) at a fraction of the
computational cost.

\subsection{Connection to Organisational and Social Theories}

\paragraph{Ricardo (1817) --- Comparative advantage.}
Ablation A3 quantifies the comparative advantage of specialised agents.
Legal and Procurement are not globally smarter than Compliance, but have a
relative advantage in detecting their specific type of risk.
That relative advantage produces 30.8\,pp of additional HDA -- exactly the
mechanism Ricardo describes for productive specialisation \citep{ricardo1817}.

\paragraph{Granovetter (1973) --- Weak ties.}
Finance is the node with betweenness\,=\,4, the most critical ``weak tie'' in
the system.
Its hypothetical removal would cut the bridge between the operational validation
community (Requester--Procurement) and the regulatory decision community
(Legal--Compliance).
Weak ties are paradoxically the most critical for system cohesion -- exactly as
\citet{granovetter1973} predicts for social networks.

\paragraph{Mancur Olson (1965) --- Collective action and free-riding.}
The HDA\,=\,0\% in M4 and M4b, even with powerful models, illustrates Olson's
collective action problem \citep{olson1965}.
When all agents have the capacity to detect risk, none assumes the
responsibility of escalating it -- each agent assumes another will do so.
The M5 solution (explicitly designating the Compliance agent as responsible for
escalation) is exactly the institutional solution Olson proposes: a specific
agent responsible for the collective good, although with \texttt{llama-3.1-8b}
the model capacity remains the limiting factor.

\subsection{Limitations and Unmet Objectives}

The most relevant limitation is that none of the evaluated configurations
surpassed the target HDA\,$>$\,90\%.
The maximum HDA achieved was 30.8\%, with M2.
Notably, M4 -- which uses a more powerful decision model -- reached only
HDA\,=\,5.1\%.
This apparently paradoxical result points to a research hypothesis with broader
implications: role specialisation in multi-agent architectures may inhibit the
detection of systemic risk.
When each agent is responsible only for its domain, cases requiring
cross-cutting evaluation -- GDPR combined with contractual risk and high amount,
for example -- are not escalated because no specialised agent perceives the
complete signal.
The single agent, by integrating all context, detects the combination of
factors even without an explicit label for it.
This connects directly with Olson's collective action problem
\citeyearpar{olson1965}: when responsibility is distributed, no one assumes it.

\begin{quote}
\textit{Hypothesis for future work: Reliable detection of systemic risk in
multi-agent architectures requires an explicit cross-cutting synthesis mechanism
-- not a more powerful agent, but an information flow that preserves the overall
picture before the escalation decision.}
\end{quote}

The uncertainty classifier M5b, which achieved HDA\,=\,23.1\% with 0 false
positives, is a first step in that direction.
Its extension with specific regulatory factors (GDPR, AESA, CE certification)
constitutes the most promising line of work to approach the HDA\,$>$\,90\%
target.

The experiment covers only the procurement process.
The generalisation question -- whether findings on compute equivalence and model
diversity are specific to procurement or hold in processes with different
decision structures -- remains open.
Processes with greater parallelism could show more pronounced multi-agent
advantages with small models, given that context fragmentation would harm less
when subproblems are independent.

The variability between M2 runs ($\sigma$\,ARR\,=\,11.3\,pp) with
\texttt{temperature=0.0} reveals that single-run results with 8B models must be
interpreted with caution.
M4 is the only LLM system with acceptable variability for production
environments ($\sigma$\,ARR\,=\,0.9\,pp).
Finally, the six enterprise tools were implemented as mock APIs with
deterministic logic, which guarantees experiment reproducibility but introduces
a gap with production reality -- real ERP/CRM systems present variable
latencies, heterogeneous data, and inconsistent states not represented in the
evaluated dataset.

%% ─────────────────────────────────────────────────────────────────────────────
%% 10. CONCLUSIONS AND FUTURE WORK
%% ─────────────────────────────────────────────────────────────────────────────
\section{Conclusions and Future Work}
\label{sec:conclusions}

This work started from a concrete question: can a system of LLM agents resolve
in seconds what today takes days in an enterprise procurement process?
The answer is affirmative, but the nuances are more instructive than the
confirmation.

Model M4 -- five specialised agents with \texttt{llama-3.1-8b} in auxiliary
roles and \texttt{llama-3.3-70b} as the final decision-maker -- achieved
ARR\,=\,98.7\% with a variability of $\pm$0.9\,pp across three independent
replicas.
It processes requests in 1.74 seconds at an estimated cost of \$0.0065 per
case, improving the error rate of the RPA baseline by 8\,pp.
The system surpasses the rule-based model in decision quality without
sacrificing processing speed.

The most unexpected finding was the behaviour of M3.
The working hypothesis was that five specialised agents would outperform a
single agent -- and that was not the case.
With \texttt{llama-3.1-8b}, context fragmentation harmed decision quality more
than specialisation improved it.
This forced a reframing of the underlying question: not whether multi-agent
architecture is better in absolute terms, but above \emph{what model capacity
threshold} it is.
The data suggest that threshold lies between 8B and 70B parameters for this
type of task.

The second surprising finding was the HDA pattern.
M2 -- single agent, small model -- detected more escalation cases than M4,
which has a significantly more powerful decision model.
After reviewing the cases, the emerging explanation was: the single agent
integrates all risk factors in a single context, while in the multi-agent
architecture each agent only perceives its fragment.
Cases requiring escalation are precisely those presenting cross-cutting risk --
no specialised agent detects it in its entirety.
This reframes the HDA problem: it is not a prompt engineering question, but an
information architecture question.

\subsection{Contributions}

The four main contributions of this work are:
\begin{enumerate}[label=(\arabic*)]
  \item a complete multi-agent framework for procurement with five specialised
        agents, six simulated enterprise tools, and HiL support, implemented in
        CrewAI with Python 3.12;
  \item a dataset of 50 cases with annotated ground truth across three
        complexity levels, including regulatory exceptions (GDPR, AESA, CE
        certification) not present in generic benchmarks;
  \item empirical evidence on compute equivalence, model diversity, and role
        specialisation via an executed ablation study; and
  \item statistical robustness analysis with three independent replicas that
        quantifies the instability of 8B models ($\sigma$\,ARR\,=\,11.3\,pp)
        versus the stability of the mixed architecture with a 70B decision-maker
        ($\sigma$\,ARR\,=\,0.9\,pp).
\end{enumerate}

\subsection{Methodological Reflection}

If the experiment were repeated, three changes would be prioritised.
First, design the HiL escalation mechanism as part of the information
architecture from the outset -- not as a prompt add-on.
The M5b classifier, built in response to M5 results, should have been the
starting point of the design.
Second, start with the most powerful model in all agents to establish the
performance ceiling, then progressively reduce to lightweight models, rather
than the reverse order followed.
Third, include statistical replicas from the first run: the variability observed
in M2 with \texttt{temperature=0} was not expected, and detecting it late
required reinterpreting results that initially appeared stable.

\subsection{Future Work}

Priority research directions are as follows.
First, evaluate whether HDA\,$>$\,90\% is achievable with the current
architecture using powerful models in all agents, or whether the problem is
structural -- i.e., inherent to context distribution -- regardless of model.
Second, extend the M5b uncertainty classifier by incorporating specific
regulatory factors as explicit escalation signals, addressing the gap identified
in the 10 undetected HiL cases.
Third, replicate the experiment in four additional business processes (onboarding,
incident management, compliance check, and financial analysis) to determine
whether findings on compute equivalence and model diversity generalise beyond
procurement.
Fourth, connect the system with live ERP and CRM APIs to evaluate behaviour
with ambiguous, incomplete, or inconsistent data.
Fifth, investigate architectures with explicit cross-cutting synthesis -- an
integrator agent that receives the outputs of all specialised agents before the
escalation decision -- as a way to preserve both specialisation and systemic
risk awareness.

\begin{quote}
\textbf{Conclusion}: Multi-agent architecture is not superior to a single agent
by default.
It is superior when the underlying model has sufficient capacity to benefit from
specialisation.
With powerful models, M4 combines the speed of RPA with the reasoning capability
of LLMs, resolving in seconds processes that today require days.
The path to HDA\,$>$\,90\% does not lie in improving the models that inhabit the
agents, but in redesigning how information flows between them.
\end{quote}

%% ─────────────────────────────────────────────────────────────────────────────
%% REFERENCES
%% ─────────────────────────────────────────────────────────────────────────────
\bibliographystyle{elsarticle-harv}
\bibliography{projener_references}

%% ─────────────────────────────────────────────────────────────────────────────
%% APPENDICES
%% ─────────────────────────────────────────────────────────────────────────────
\appendix

\section{Agent System Prompts}
\label{app:prompts}

The following system prompts were used for each agent.
These are provided as required by Q1 venue standards for reproducibility.

\paragraph{Requester Agent.}
\begin{lstlisting}[language={}]
role: "Gestor de Solicitudes"
goal: "Registrar y estructurar solicitudes de compra."
backstory: "Eres el primer punto de contacto. Registras
           solicitudes y las preparas para el equipo."
\end{lstlisting}

\paragraph{Procurement Agent.}
\begin{lstlisting}[language={}]
role: "Especialista en Proveedores"
goal: "Validar proveedores y buscar alternativas."
backstory: "Proveedor: estado=X, opera=Y"
           [API catalogue result]
\end{lstlisting}

\paragraph{Finance Agent.}
\begin{lstlisting}[language={}]
role: "Analista Financiero"
goal: "Verificar presupuesto y aprobar o proponer reasignacion."
backstory: "Presupuesto: disponible=X EUR, aprobado=Y"
           [budget API result]
\end{lstlisting}

\paragraph{Legal Agent.}
\begin{lstlisting}[language={}]
role: "Asesor Legal"
goal: "Revisar clausulas contractuales y detectar riesgos."
backstory: "Legal: riesgo=X, firmar=Y"
           [contract validation API result]
\end{lstlisting}

\paragraph{Compliance Agent (decision-maker, M3/M5).}
\begin{lstlisting}[language={}]
role: "Auditor de Compliance"
goal: "Verificar normativa y decidir."
backstory: "Normativa: cumple=X, riesgo=Y.
            [M5: ESCALA si conflicto interes, GDPR,
             importe >50k EUR, riesgo critico]"
output: JSON {decision, razon, escala_hil}
\end{lstlisting}

\section{Reproducibility Parameters}
\label{app:repro}

\begin{itemize}
  \item Primary model: \texttt{groq/llama-3.1-8b-instant}
  \item Decision model M4/M4b: \texttt{groq/llama-3.3-70b-versatile}
  \item Temperature: 0.0 (initial verification: 5/5 cases identical in single
        session: C01, C21, C39, C43, C50)
  \item Temperature note: single-session verification does not capture
        inter-session variability.
        The statistical robustness analysis (Section~\ref{sec:setup}) reveals
        $\sigma$\,ARR\,=\,11.3\,pp for \texttt{llama-3.1-8b-instant} and
        $\sigma$\,ARR\,=\,0.9\,pp for \texttt{llama-3.3-70b-versatile}.
        \texttt{llama-3.3-70b} is recommended for production environments.
  \item Inter-case pause: M1--M2x: 4\,s; M3--M5: 5--8\,s; M4b: 8\,s
  \item Maximum retries per case: 3
  \item Total LLM calls (single run): M1=0; M2=50; M2x=50; M3=250; M3b=150;
        M4=250; M4b=250; M5=250; M5b=100
  \item Total LLM calls (statistical analysis): M1$\times$3=0;
        M2$\times$3=150; M4$\times$3=750 additional replica calls
  \item Variants executed: 9 models $\times$ 50 cases\,=\,450 total cases
        + 200 cases in statistical replicas
  \item Complete code: \url{https://github.com/Edurne0781/multiagent-procurement-llm}
        (Jupyter Lab, Python 3.12)
\end{itemize}

\end{document}
